\section{Розділ~\ref{sec:sleep}. Сон}

Режими сну, повтори й зменшення синапсів виміряно записами нейронів, мікродіалізом і електронною мікроскопією; розклад сну й повільні гойдалки описано моделями книги, а призначення сну --- здебільшого гіпотези.

{\footnotesize
\setlength{\tabcolsep}{3pt}\renewcommand{\arraystretch}{1.15}
\begin{longtable}{>{\raggedright}p{0.5cm}>{\raggedright}p{4.0cm}>{\raggedright}p{5.6cm}>{\raggedright}p{2.3cm}>{\raggedright\arraybackslash}p{1.7cm}}
\toprule
№ & Твердження & Дослід і що виміряно & Джерело & Статус \\
\midrule\endfirsthead
\toprule
№ & Твердження & Дослід і що виміряно & Джерело & Статус \\
\midrule\endhead
1 & Уві сні нейрони кори не мовчать; змінюється режим роботи & Тривалі записи нейронів кори в щурів: у повільному сні частота лишається ненульовою, періоди роботи чергуються з періодами тиші; після тривалого неспання частота вища в усіх станах, після сну --- нижча & \cite{Vyazovskiy2009} & виміряно \\
2 & \nt{ACET} високий удень і у швидкому сні, низький у повільному (таблиця~\ref{tab:sleepmodes}) & Мікродіаліз у корі вільно рухливих кішок: під час неспання викид ацетилхоліну на $100$--$175\%$ вищий, ніж у повільному сні, а у швидкому сні --- приблизно на рівні спокійного неспання & \cite{Marrosu1995} & виміряно \\
3 & \nt{NORA} і \nt{SERO} затихають у повільному сні й майже мовчать у швидкому & Запис окремих \nt{NORA}-нейронів у щурів і \nt{SERO}-нейронів у кішок за стадіями сну: частота максимальна під час неспання, нижча в повільному сні й майже нульова у швидкому & \cite{AstonJonesBloom1981,McGintyHarper1976} & виміряно \\
4 & У швидкому сні м'язи вимкнено гальмуванням через \nt{GABA} і \nt{GLYC} & У щурів вимірювали активність \nt{ACET}-нейронів жувального м'яза й подавали блокатори просто до них: параліч знімається, лише якщо заблоковано і \nt{GABA}-рецептори обох типів, і \nt{GLYC}-рецептори & \cite{BrooksPeever2012} & виміряно \\
5 & Тиск сну $S$ --- конденсатор із двома порогами~\eqref{eq:sleeppressure}, $\tau_r=18.2$~год, $\tau_d=4.2$~год & Модель двох процесів; сталі часу отримано з того, як потужність повільних хвиль ЕЕГ зростає з неспанням і падає уві сні, а форму порогів --- з часу спонтанного пробудження & \cite{Borbely1982,Daan1984} & теорія \\
6 & Машина сама виходить на $16.1$~год неспання і $8.0$~год сну (рис.~\ref{fig:twoprocess}) & Розрахунок за~\eqref{eq:sleeppressure} з гойдальними порогами, скрипт \texttt{models/sleep\_models.py}: $16.1$~год неспання і $7.9$--$8.0$~год сну & --- & модель \\
7 & Після $40$~год без сну $S=0.90$, але сон триває лише $8.8$~год: борг повертається глибиною, а не тривалістю & Модель: \texttt{models/sleep\_models.py} дає $S=0.90$ і $8.8$~год. Дослід: після $40.5$~год без сну у восьми людей у першу відновлювальну ніч потужність повільних хвиль ЕЕГ значуще вища за звичайну & \cite{Borbely1981} & модель \\
8 & Фізично $S$ --- аденозин & Мікродіаліз у кішок: позаклітинний аденозин в одній із ділянок, що керують неспанням, зростає за час неспання, зростає далі під час позбавлення сну й падає під час сну. Що $S$ --- лише аденозин, не показано & \cite{PorkkaHeiskanen1997,Dunwiddie2001} & гіпотеза \\
9 & У повільному сні кора гойдається: робота кілька сотень мс --- тиша, рідше ніж раз на секунду ($<1$~Гц, під наркозом частіше $0.3$--$0.4$~Гц) & Внутрішньоклітинний запис у кішок під наркозом: деполяризація $0.8$--$1.5$~с і тривала гіперполяризація, ритм $<1$~Гц, найчастіше $0.3$--$0.4$~Гц; те саме чергування роботи й тиші видно в природному сні (рядок~1) & \cite{Steriade1993,Vyazovskiy2009} & виміряно \\
10 & Гойдалки --- група зі зворотним зв'язком і втомою~\eqref{eq:updown}; за $b=5$ період $1.1$~с (значення моделі) і робота близько $35\%$ часу, за $b=1$ робота рівна (рис.~\ref{fig:updown}) & Розрахунок, скрипт \texttt{models/sleep\_models.py}: період $1.11$~с, частка часу роботи $0.35$ (середнє за цілим числом періодів); за $b=1$ частка роботи $1.0$. Період моделі коротший за виміряний (рядок~9) & --- & модель \\
11 & \nt{ACET} вимикає гойдалки й переводить кору в рівну роботу & Стимуляція ядра \nt{ACET}-нейронів у кішки блокувала повільний ритм у $79\%$ нейронів кори й переводила їх у рівну роботу, переважно через зникнення тривалих гіперполяризацій. Ацетилхолін пригнічує повільні калієві струми, що створюють утому & \cite{SteriadeAmzica1993,McCormick1992} & виміряно \\
12 & Спалахи ритму $7$--$15$~Гц породжує петля \nt{GLUT}--\nt{GABA} між корою й релейним вузлом & У зрізах релейного вузла зору в тхора спалахи створюються відповідними пачками релейних клітин на гальмування від сусіднього \nt{GABA}-ядра, яке вони самі збуджують & \cite{vonKrosigk1993,SteriadeMcCormickSejnowski1993} & виміряно \\
13 & Повтори дня йдуть із перемотуванням: у гіпокампі приблизно в $20$ разів швидше, у корі --- в $6$--$7$ разів & У повільному сні повторювані послідовності нейронів гіпокампа йдуть стисненими в часі. Після бігу доріжкою послідовності нейронів місця повторювалися в тому самому порядку короткими спалахами по $\sim100$~мс, стисненими приблизно в $20$ разів; у корі стиснення в $6$--$7$ разів & \cite{Nadasdy1999,LeeWilson2002,Euston2007} & виміряно \\
14 & Посилення узгодженості повторів із корою поліпшує пам'ять (причиновий зв'язок) & У щурів після навчання стимуляцією, прив'язаною до повторів, посилювали їхній зв'язок із повільними хвилями й спалахами ритму в корі: наступного дня пам'ять була високою, у контрольних --- на рівні випадку & \cite{Maingret2016} & виміряно \\
15 & Призначення повторів --- перемішане навчання повільної пам'яті & Теорія двох систем навчання й розрахунки на штучних мережах: послідовне навчання псує старе, перемішане --- ні. Прямого досліду на мозку, що повтори потрібні саме для цього, немає & \cite{McClelland1995} & гіпотеза \\
16 & Після сну синапси зменшуються пропорційно розміру, великі майже не змінюються & Тривимірна електронна мікроскопія $6920$ синапсів кори миші: площа контакту після сну на $\sim18\%$ менша, ніж після неспання, зменшення пропорційне розміру, великі стійкі синапси не змінюються & \cite{deVivo2017} & виміряно \\
17 & Головне завдання сну --- перенормування ваг & Гіпотеза синаптичного гомеостазу; рядки 1 і 16 з нею узгоджуються, але не доводять, що це головна функція & \cite{TononiCirelli2014} & гіпотеза \\
18 & Швидкий сон --- випробування моделі світу на стенді & Режим (таблиця~\ref{tab:sleepmodes}) виміряно; що мережа проганяє модель світу, --- теоретичне припущення, досліду немає & \cite{Hobson2009} & гіпотеза \\
19 & Промивання мозку уві сні: питання відкрите & Один дослід: уві сні міжклітинний простір на $60\%$ ширший і очищення швидше. Інший, з іншим способом вимірювання: уві сні й під наркозом очищення помітно повільніше. Результати суперечать один одному & \cite{Xie2013,Miao2024} & гіпотеза \\
20 & Місцевий сон: ділянки кори тварини, що не спить, ненадовго провалюються в тишу & У щурів після тривалого неспання нейрони окремої ділянки кори коротко замовкають, як у повільному сні, з повільною хвилею в місцевій ЕЕГ; у такі моменти щур частіше помиляється в задачі & \cite{Vyazovskiy2011} & виміряно \\
\bottomrule
\end{longtable}}
